% Part: sets-functions-relations
% Chapter: relations
% Section: relations-as-sets

\documentclass[../../../include/open-logic-section]{subfiles}

\begin{document}

\olfileid{sfr}{rel}{set}
\olsection{اړيکې د سټونو په توګه}

\begin{explain}
په \olref[sfr][set][imp]{sec} کښې مو د ځينو مهمو سټونو يادونه کړې وه:
$\Nat$، $\Int$، $\Rat$، $\Real$۔ بې شکه به تاسو ته د دې سټونو
د ځينو د \psOblique{!!{element}s} ترمنځ په زړه پورې اړيکې يادې وي۔
د بېلګې په توګه، په دې هر سټ باندې يوه پوره معياري \emph{ترتيبي اړيکه}
شته۔ د \emph{هماغه شے کېدو} اړيکه هم شته چې هر څيز ئې له خپل ځان
سره لري او له بل هېڅ څيز سره ئې نۀ لري۔ مونږ به له ډېرو نورو په زړه
پورو اړيکو سره هم مخامخ شو، او ممکنې اړيکې تر دې هم ډېرې دي۔ خو تر
کتنې وړاندې به په دې خبره پيل وکړو چې اړيکې د سټ يو خاص ډول ګڼلے شو۔

د دې دپاره له \olref[sfr][set][pai]{sec} څخه دوه خبرې رايادې کړئ۔
لومړے، د \emph{مرتبې جوړې} مفهوم: که $a$ او $b$ راکړل شي، نو
$\tuple{a, b}$ جوړولے شو۔ دلته د غړو ترتيب \emph{خامخا} اهميت لري۔
نو که $a \neq b$ وي، $\tuple{a, b} \neq \tuple{b, a}$۔
(دا له نامرتبو جوړو، يعنې له $2$-غړيزو سټونو، سره پرتله کړئ؛ هلته
$\{a, b\}=\{b, a\}$۔) دويم، د \emph{کارتيزي حاصل ضرب} مفهوم
راياد کړئ: که $A$ او $B$ سټونه وي، نو $A \times B$ جوړولے شو۔
دا د ټولو هغو جوړو $\tuple{x, y}$ سټ دے چې $x \in A$ او $y \in B$
وي۔ په تېره بيا، $A^{2}= A \times A$ د $A$ د ټولو مرتبو جوړو سټ دے۔

اوس به په يو سټ باندې يوه خاصه اړيکه وګورو: د طبيعي عددونو په سټ
$\Nat$ باندې د $<$ اړيکه۔ د عددونو د ټولو هغو جوړو $\tuple{n, m}$
سټ په پام کښې ونيسئ چې $n<m$ وي، يعنې:
\[
  R=\Setabs{\tuple{n, m}}{n, m \in \Nat \text{ او } n<m}.
\]
د $n$ له $m$ څخه د کوچني کېدو او په $R$ کښې د جوړې $\tuple{n, m}$
د غړي کېدو ترمنځ نژدې تړاو شته، يعنې:
\[
      n<m\text{ هله او يوازې هله چې }\tuple{n, m} \in R.
\]
په رښتيا، د معلوماتو له هېڅ ضايع کېدو پرته، سټ $R$ په $\Nat$ باندې
د $<$ اړيکه \emph{ګڼلے} شو۔

په همدې طريقه، د عددونو ترمنځ د هرې اړيکې دپاره د $\Nat^{2}$ يو
فرعي سټ جوړولے شو۔ برعکس، که د عددونو د جوړو هر سټ
$S \subseteq \Nat^{2}$ راکړل شي، نو د عددونو ترمنځ ورسره سمه يوه
اړيکه شته: $n$ له $m$ سره دا اړيکه هله او يوازې هله لري چې
$\tuple{n, m} \in S$ وي۔ دا خبره لاندې تعريف ته جواز ورکوي:
\end{explain}

\begin{defn}[دوه ځاييزه اړيکه]
په سټ $A$ باندې يوه \emph{دوه ځاييزه اړيکه} د $A^{2}$ يو فرعي سټ
دے۔ که $R \subseteq A^{2}$ په $A$ باندې يوه دوه ځاييزه اړيکه وي او
$x, y \in A$ وي، نو کله کله د $\tuple{x, y} \in R$ پر ځاے
$Rxy$ (يا $xRy$) ليکو۔
\end{defn}

\begin{ex}
  \ollabel{relations}
د طبيعي عددونو د جوړو سټ $\Nat^{2}$ په يوه دوه بعدي جدول کښې داسې
ليکلے شو:
\[
  \begin{array}{ccccc}
  \mathbf{\tuple{ 0,0 }} & \tuple{ 0,1 } &
    \tuple{ 0,2 } & \tuple{ 0,3 } & \ldots\\
  \tuple{ 1,0 } & \mathbf{\tuple{ 1,1 }} &
    \tuple{ 1,2 } & \tuple{ 1,3 } & \ldots\\
  \tuple{ 2,0 } & \tuple{ 2,1 } &
    \mathbf{\tuple{ 2,2 }} & \tuple{ 2,3 } & \ldots\\
  \tuple{ 3,0 } & \tuple{ 3,1 } & \tuple{ 3,2 } &
    \mathbf{\tuple{ 3,3 }} & \ldots\\
  \vdots & \vdots & \vdots & \vdots & \mathbf{\ddots}
  \end{array}
\]
دلته مو قطر په پنډ ليک ښودلے، ځکه د $\Nat^2$ هغه فرعي سټ چې پر
قطر پرتې جوړې لري، يعنې:
\[
  \{\tuple{0,0 }, \tuple{ 1,1 }, \tuple{ 2,2 }, \dots\},
  \]
په $\Nat$ باندې \emph{د عينيت اړيکه} ده۔ (د عينيت اړيکه ډېره کارېږي؛
راځئ د هر سټ $A$ دپاره
$\Id{A}=\Setabs{\tuple{ x,x }}{x \in A}$ تعريف کړو۔) د قطر له پاسه
د ټولو پرتو جوړو فرعي سټ، يعنې:
\[
  L = \{\tuple{ 0,1 },\tuple{ 0,2 },\ldots,\tuple{ 1,2 },
  \tuple{ 1,3 }, \dots, \tuple{ 2,3 }, \tuple{ 2,4 },\ldots\},
\]
د \emph{کوچني کېدو} اړيکه ده؛ يعنې $Lnm$ هله او يوازې هله چې
$n<m$ وي۔ د قطر لاندې د پرتو جوړو فرعي سټ، يعنې:
\[
  G=\{ \tuple{ 1,0 },\tuple{ 2,0 },\tuple{
    2,1 }, \tuple{ 3,0 },\tuple{ 3,1 },\tuple{ 3,2 }, \dots\},
\]
د \emph{لويېدو} اړيکه ده؛ يعنې $Gnm$ هله او يوازې هله چې $n>m$
وي۔ د $L$ او $I$ اتحاد، چې $K=L\cup I$ ئې بللے شو، د
\emph{کوچني يا برابر کېدو} اړيکه ده: $Knm$ هله او يوازې هله چې
$n \le m$ وي۔ همدارنګه، $H=G \cup I$ د \emph{لوي يا برابر کېدو
اړيکه} ده۔ دا اړيکې، $L$، $G$، $K$ او $H$، د اړيکو خاص ډولونه دي
چې \emph{ترتيبونه} بللے شي۔ $L$ او $G$ دا خاصيت لري چې هېڅ عدد له
خپل ځان سره $L$ يا $G$ اړيکه نۀ لري (يعنې، د هر $n$ دپاره، نۀ $Lnn$
شته او نۀ $Gnn$)۔ د دې خاصيت لرونکې اړيکې \emph{نفي انعکاسي}
بللے شي، او که ترتيبونه هم وي، نو \emph{سخت ترتيبونه} بللے شي۔
\end{ex}

\begin{explain}
که څه هم ترتيبونه او عينيت مهمې او طبيعي اړيکې دي، بايد ټينګار وشي
چې زمونږ د تعريف له مخې د $A^{2}$ \emph{هر} فرعي سټ په $A$ باندې
يوه اړيکه ده، که هر څومره غيرطبيعي يا په تکلف جوړه شوې ښکاري۔
په تېره بيا، $\emptyset$ په هر سټ باندې يوه اړيکه ده (\emph{تشه اړيکه}،
چې د غړو هېڅ جوړه ئې نۀ لري)۔ پخپله $A^{2}$ هم په $A$ باندې يوه
اړيکه ده (چې هره جوړه ئې لري)، او \emph{ټولشموله اړيکه} بللے شي۔
خو د دې په شان يو څيز هم اړيکه شمېرلے شي:
$E=\Setabs{\tuple{n, m}}{n>5 \text{ يا } m \times n \ge 34}$۔
\end{explain}

\begin{prob}
  په سټ $\Pow{\{a, b, c\}}$ باندې د $\subseteq$ اړيکې
  !!{element}s په لړ کښې وليکئ۔
\end{prob}

\end{document}
